\subsection{Cauchy 滤子}

一旦集合 X 赋有了一致结构，我们就可以定义（相对于该结构的）X 的“小”子集指的是什么：X 的“小”子集是指其中所有点彼此“非常接近”的那种子集。精确地说：

定义 1 若 X 是一致空间而 V 是 X 的一致邻域，则称 X 的子集 A 是 V-小集，如果 A 的任意两点都 V-邻近（换言之，A × A ⊂ V）。

命题 1 在一致空间 X 中，若两个集合 A 与 B 都是 V-小集且相交，则其并 A ∪ B 是 V-小集。

设 \(x\) 与 \(y\) 是 \(A \cup B\) 的任意两点，并设 \(Z \in A \cap B\)。则 \((x, z) \in V\) 且 \((z, y) \in V\)，从而 \((x, y) \in \mathring{V}\)。

定义 2 滤子 \(\mathfrak{F}\)（一致空间 \(\mathbf{X}\) 上的）称为 Cauchy 滤子，如果对 X 的每个一致邻域 V，存在 X 的一个子集，它是 V-小集且属于 \(\mathfrak{F}\)。

这里我们同样可以借助“足够小的集合”与“要多小有多小的集合”这些说法使语言更形象；于是定义 2 可以改述为：Cauchy 滤子就是包含任意小的集合的滤子。

无穷序列 \((u_n)\)（一致空间 \(X\) 的点列）称为 Cauchy 序列，如果与该序列相伴的初等滤子是 Cauchy 滤子。这等价于说：对每个一致邻域 \(V\)（\(X\) 的），存在整数 \(n_0\)，使得对一切满足 \(m \geqslant n_0\) 与 \(n \geqslant n_0\) 的整数 m 与 n，都有 \((u_m, u_n) \in V\)。

命题 2 在一致空间 X 上，每个收敛滤子都是 Cauchy 滤子。

若 \(x\) 是 \(X\) 的任意一点，且 \(V\) 是 \(X\) 的任意对称一致邻域，则邻域 \(V(x)\)（\(x\) 的）是 \(V\)-小集。若 \(\mathfrak{F}\) 是收敛于 \(x\) 的滤子，则 \(\mathfrak{F}\) 中有一个集合含于 \(V(x)\)，因而是 \(V\)-小集。

显然，每个细于某个 Cauchy 滤子的滤子都是 Cauchy 滤子。

命题 3 设 \(f: \mathbf{X} \to \mathbf{X}'\) 是一致连续映射。则在 \(f\) 之下，\(\mathbf{X}\) 上 Cauchy 滤子基的像是 \(\mathbf{X}'\) 上的 Cauchy 滤子基。

设 \(g = f \times f\)。若 \(V'\) 是 \(X'\) 的一致邻域，则 \(\overline{g}^1(V')\) 是 \(X\) 的一致邻域，且在 \(f\) 之下，\(\overline{g}^1(V')\)-小集的像是 \(V'\)-小集；由此即得结论。

特别由此可得：若把一致空间 X 的一致结构换成较粗的一致结构，则关于原来一致结构的每个 Cauchy 滤子关于新一致结构仍是 Cauchy 滤子。

这一事实很容易按下述形式记住：一致结构越细，Cauchy 滤子就越少。

命题 4 设 X 是一个集合，\((\mathbf{Y}_i)_{i\in I}\) 是一族一致空间，且对每个 \(i\in I\)，\(f_i\) 是 X 到 \(\mathbf{Y}_i\) 的映射。设 X 赋有最粗的一致结构 \(\mathcal{U}\)，使诸 \(f_i\) 都一致连续。则 X 上的滤子基 \(\mathcal{B}\) 是 Cauchy 滤子基的充要条件是：\(f_i(\mathcal{B})\) 是 \(\mathbf{Y}_i\) 上的 Cauchy 滤子基（对每个 \(i\in I\)）。

由命题 3，条件是必要的。反之，设条件满足，并设 \(\mathbf{U}(\mathbf{V}_{\iota_1},\ldots ,\mathbf{V}_{\iota_n})\) 是一致结构 \(\mathfrak{U}\) 的一个一致邻域{[}§ 2 第 3 目公式 (1){]}。由假设，对每个指标 \(k\)，存在集合 \(\mathbf{M}_k\in \mathfrak{B}\)，使 \(f_{\iota_k}(\mathbf{M}_k)\) 是 \(\mathbf{V}_{\iota_k}\)-小集 ( \(\iota \leqslant k\leqslant n\) )。设 \(\mathbf{M}\) 是 \(\mathfrak{B}\) 中含于诸 \(\mathbf{M}_k\)（\(\iota \leqslant k\leqslant n\)）的一个集合；则对每对点 \(x,x^{\prime}\)（\(\mathbf{M}\) 的），有 \([f_{\iota_k}(x),f_{\iota_k}(x')]\in \mathbf{V}_{\iota_k}\)（\(\iota \leqslant k\leqslant n\)），从而

\[
(x, x ^ {\prime}) \in \mathrm{U} (\mathrm{V} _ {t _ {1}}, \dots , \mathrm{V} _ {t _ {n}}).
\]

证明完毕。

推论 1 若一致空间 X 上的 Cauchy 滤子在 X 的子集 A 上诱导出一个滤子，则该滤子是一致子空间 A 上的 Cauchy 滤子。

推论 2 滤子基 \(\mathfrak{B}\)（一致空间的积 \(\prod_{i\in I}X_i\) 上的）是 Cauchy 滤子基的充要条件是：对每个 \(i\in I\)，\(\operatorname{pr}_i(\mathfrak{B})\) 是 \(X_i\) 上的 Cauchy 滤子基。

